\section{Other methods: tensor networks, CLIK, eigenvalues without LAPACK}
\label{app:other}

These methods were developed to go beyond the exact sizes. None has done so yet: the dissipated scar state needs nearly the full bond dimension as a matrix product state (Sec.~\ref{sec:tn-results}), CLIK has been run only at sizes where the full Liouvillian spectrum is known, and the forward/backward code has only its $L = 6$ validation. They are kept here with their tests.

\subsection{Tensor networks: $\rho$ as a matrix product state, TDVP, influence matrix, CLIK}
\label{sec:tn}

\paragraph{$\rho$ as an MPS.} The Pauli coordinates factorize over sites, $r_P = r_{p_0p_1\dots p_{L-1}}$ with $p_j \in \{0, x, y, z\}$ the factor of $P$ on site $j$, so $\rho$ is a vector in $(\mathbb R^4)^{\otimes L}$ and can be written as a matrix product state
\begin{equation}
r_{p_0\dots p_{L-1}} = A_0^{p_0}A_1^{p_1}\cdots A_{L-1}^{p_{L-1}},
\end{equation}
with $\chi_{j}\times\chi_{j+1}$ matrices $A_j^{p}$ ($\chi_0 = \chi_L = 1$; ITensor, site type \texttt{PauliVec}, \texttt{julia/lindblad\_tdvp.jl}). The singular values of the matrix $r_{(p_0\dots p_{j-1}),(p_j\dots p_{L-1})}$ are the operator Schmidt values of $\rho$ across the cut between sites $j-1$ and $j$ (Appendix~\ref{sec:methods}, the strings being orthogonal with equal norms). The representation is exact for $\chi_j = 4^{\min(j, L-j)}$. A truncation to bond dimension $\chi$ keeps the $\chi$ largest Schmidt values and, in the normalized vector, loses the weight $\varepsilon^2 = \sum_{\ell>\chi}w_\ell$. The operator entanglement $S_{\mathrm{op}}$ of $\rho(t)$ therefore sets the cost: $\chi \gtrsim e^{S_{\mathrm{op}}}$ is needed for small errors. Expectation values are contractions: $\langle\sigma^\alpha_j\rangle = r_{0\dots\alpha_j\dots0}$, read off by projecting every other site on $p = 0$.

\paragraph{The generator as an MPO.} Left multiplication by a Pauli matrix permutes Pauli matrices with phases, $\sigma^a\sigma^b = \sum_c (A^a + iB^a)_{cb}\,\sigma^c$ with real $4\times4$ matrices $A^a$, $B^a$ ($A^a$ the symmetric part: $\sigma^a\sigma^0 = \sigma^a$, $\sigma^a\sigma^a = \sigma^0$; $B^a$ the antisymmetric part from $\sigma^a\sigma^b = i\epsilon_{abc}\sigma^c$), and right multiplication gives $A^a - iB^a$. For a string $P$ with factors $a_j$, $[P, X]$ acts on the coordinates as $\bigotimes_j(A^{a_j} + iB^{a_j}) - \bigotimes_j(A^{a_j} - iB^{a_j})$, so
\begin{equation}
-i[P,\cdot\,] \;\longrightarrow\; 2\operatorname{Im}\bigotimes_j\big(A^{a_j} + iB^{a_j}\big),
\end{equation}
a real operator, and the dephasing is the diagonal $-2\gamma\sum_j(\text{projector on } p_j \in \{x, y\})$. Summing over the strings of $H$ gives $G$ as a real matrix product operator; the periodic boundary enters as long-range terms.

\paragraph{TDVP: forward and backward sweeps.} The time-dependent variational principle (Dirac--Frenkel) keeps the state on the manifold $\mathcal M_\chi$ of MPS of bond dimension $\chi$ by projecting the generator onto its tangent space,
\begin{equation}
\frac{d}{dt}\ket{r} = \mathcal P_{T_r\mathcal M_\chi}\,G\ket{r}.
\end{equation}
In mixed-canonical gauge the tangent-space projector decomposes into one-site and bond terms, $\mathcal P_T = \sum_{j}\mathcal P^{(1)}_j - \sum_{j}\mathcal P^{(0)}_{j,j+1}$, and the projector-splitting integrator solves the terms one after the other. In a left-to-right sweep, the site tensor $A_j$ is evolved \emph{forward} by $\delta t$ with its local effective generator, split by a QR decomposition into a left-orthonormal tensor and a bond matrix, and the bond matrix is evolved \emph{backward} by $\delta t$ (the minus sign of $\mathcal P^{(0)}$) before it is absorbed into $A_{j+1}$. A right-to-left sweep with $\delta t/2$ after a left-to-right one with $\delta t/2$ makes the step symmetric and second order. The two-site variant evolves two-site tensors forward and single-site tensors backward, and splits the two-site tensor by an SVD, so the bond dimension can grow up to $\chi_{\max}$; the truncation is the only approximation beyond the projection. The local exponentials are applied with Krylov (Arnoldi, $G$ is not symmetric) methods. Two points matter here:
\begin{itemize}
\item \emph{Projection error at small bonds:} terms of $G$ that act on more than two sites (the range-4 terms of $H$, the long-range periodic terms) leave the two-site tangent space while the bonds are small, and two-site TDVP then misses them. A global subspace expansion (enlarging the bond basis with Krylov vectors $G^k\ket r$ of the full generator before the sweep) removes this; once every bond has its full dimension TDVP is exact.
\item \emph{Truncation error:} at fixed $\chi$ it is set by the operator Schmidt spectrum of $\rho(t)$, i.e.\ by $S_{\mathrm{op}}$.
\end{itemize}
Backward (Heisenberg) evolution of an observable uses the same machinery with $G^{\mathsf T}$, i.e.\ with $H \to -H$ (Sec.~\ref{sec:forward-backward}).

\paragraph{Influence matrix and transverse contraction.} The standard contraction follows time: $\rho$ (or $A$) is evolved step by step as an MPS in space, and its cost is set by the operator entanglement in space. The influence-matrix approach contracts the same network along space instead.
\begin{itemize}
\item \emph{The space--time network.} With a Trotter step $\delta t$, $e^{\Lv\delta t} \approx \mathcal D\,\mathcal G$, where $\mathcal G$ is a brickwork of local unitary channels $X \mapsto U_gXU_g^\dagger$ (the terms of $H$ grouped into gates on a few neighbouring sites) and $\mathcal D = \bigotimes_j e^{\D_j\delta t}$ the single-site dephasing. In the Pauli (vectorized) representation each channel is a real tensor whose legs are 4-dimensional: every leg carries both the ket and the bra index of $\rho$, i.e.\ the forward and backward branches of the Keldysh contour folded together. Then
\[ \langle A_x\rangle_{N\delta t} = \big\langle \mathds 1\big|\,A_x\,(\mathcal D\mathcal G)^N\,\big|\rho_0\big\rangle \]
is a two-dimensional network of $N$ time layers by $L$ sites, closed by the product state $\rho_0$ at the bottom and by the trace $\langle\mathds 1|$ (identity on every site except $x$) at the top.
\item \emph{Light cone.} A unitary channel acting on identities returns identities, and the dephasing does not change $\mathds 1$ (trace preservation, unitality). Gates outside the backward light cone of the site $x$ and the final time therefore cancel, and only a cone of width $\approx 2vN\delta t$ contributes.
\item \emph{Transverse contraction.} Grouping the network into columns, one per site (or unit cell), each column is a spatial transfer matrix $T_j$ acting on time-like vectors of dimension $4^N$, the full history of the boundary between two sites. For a translation-invariant infinite chain, $\langle A_x\rangle = \langle \mathrm{IM}_L|\,T_x[A]\,|\mathrm{IM}_R\rangle$, where the influence matrices $\langle\mathrm{IM}_L|$, $|\mathrm{IM}_R\rangle$ are the dominant (fixed-point) eigenvectors of $T$ from the left and from the right. They are found by applying $T$ repeatedly to a boundary vector and compressing the result as an MPS \emph{in time}: the transverse contraction. With site-dependent couplings (our random $c$) the fixed point is replaced by the product of the transfer matrices of the actual sites, applied from the boundaries.
\item \emph{Meaning and cost.} The influence matrix is the Feynman--Vernon influence functional of the rest of the chain on the chosen site: it encodes all memory effects of the many-body environment on the local dynamics. Its cost is set by its \emph{temporal entanglement} (the entanglement of the IM as an MPS across a time cut), not by the spatial entanglement of $\rho(t)$. For a site coupled to a perfectly dephasing (Markovian) environment the IM is a product in time. Strong local dephasing lowers temporal entanglement, and the approach reaches $L = \infty$ directly. That is what the finite-size cap (Sec.~\ref{sec:theory-weak}) prevents with exact methods at small $\gamma$; the influence matrix has not been implemented.
\end{itemize}

\paragraph{CLIK: eigenvalues from time-evolved states.} CLIK (Westhoff, Moroder, Schollw\"ock and Paeckel) computes slow eigenvalues of $\Lv$ without diagonalizing it. By Eq.~\eqref{eq:spectral}, the evolved states $\rho_n = e^{\Lv z_n}\rho_0$, sampled at $z_n = n\,\delta t$, contain each mode $R_k$ with weight $e^{\lambda_kz_n}$: the slow modes survive, and the span of the samples approximates the slow invariant subspace (a Krylov-like space generated by $e^{\Lv\delta t}$). On this span $\Lv$ is projected (Rayleigh--Ritz, Galerkin): writing $r = \sum_nc_n\rho_n$ and demanding $\langle\rho_m, (\Lv - \lambda)r\rangle = 0$ for all $m$ gives the generalized eigenproblem
\begin{equation}
\mathsf L\,c = \lambda\,\mathsf M\,c, \qquad \mathsf M_{mn} = \langle\rho_m, \rho_n\rangle, \qquad \mathsf L_{mn} = \langle\rho_m, \Lv\rho_n\rangle,
\end{equation}
solved after removing the near-zero eigenvalues of the Gram matrix $\mathsf M$ (the samples are nearly linearly dependent). Two evolutions are combined: a physical state ($\Tr = 1$) in real time, for the steady state and slow modes, and a traceless state along the complex contour $z = t\,e^{-i\alpha}$. Along it mode $k$ is damped at $\xi_k = -\operatorname{Re}\lambda_k\cos\alpha - \operatorname{Im}\lambda_k\sin\alpha$, so modes with $\operatorname{Im}\lambda_k > 0$ are enhanced, provided $\alpha < \alpha_{\max} = \min_k(\varphi_k - \pi/2)$ ($\lambda_k = |\lambda_k|e^{i\varphi_k}$), beyond which some mode grows. Left eigenvectors come from the same construction with $\Lv^\dagger$ (backward evolution, Sec.~\ref{sec:forward-backward}). With exact states the method is limited only by the time needed to separate the modes ($T \sim$ a few $/\Delta$, $\Delta$ the gap between the wanted modes and the rest); with MPS states the span has an effective bond dimension of order $N_K\chi$, but every state carries its truncation error, which sets the noise floor of the Gram matrix. Implementation: \texttt{julia/clik.jl} (dense backend for $L \le 6$ with exact propagators, MPS backend with two-site TDVP and global expansion), following \texttt{data/docs/CLIK\_MPS/IMPLEMENTATION.md}.


\subsection{Results of CLIK}
\label{sec:tn-results}

\paragraph{CLIK} (current coefficients, random $c$, $\gamma = 0.01$, $\delta t = 0.5$, $K = 4$ product terms per initial state; errors against the exact spectrum of $\Lv$).
\begin{itemize}
\item \textbf{Checks.} The Julia $\Lv$ (MPO, contracted) has the Rust ED spectrum to $5\times10^{-12}$ ($L = 4, 5, 6$, ring and open chain). At the full bond dimension $4^{L/2}$ the MPS backend reproduces the dense Krylov matrices to $10^{-12}$ ($L = 4$) and $10^{-13}$ ($L = 6$), but only with the global subspace expansion; without it two-site TDVP misses the range-4 terms while the bonds are small ($10^{-2}$ error from the first step).
\item \textbf{The complex contour does not help here.} From the ED, $\alpha_{\max} = 0.0020$, $0.0031$, $0.0027$ (ring, $L = 4, 5, 6$) and $0.0037$ (open, $L = 6$), set by bulk modes with $\operatorname{Im}\lambda \approx 15$--$23$, not by the slow modes. The boost of slow modes is therefore negligible and CLIK works as a real-time Krylov method, converging after $T \approx 3$--$5/\Delta \approx 200$--$300$ ($\Delta \approx 0.015$--$0.018$).
\item \textbf{Ring, random initial states, dense, $T = 400$:} $\lambda = 0$ to $10^{-15}$; the edge mode $-0.02 + i$ to $3\times10^{-14}$, $10^{-10}$, $10^{-9}$ ($L = 4, 5, 6$); the slowest bulk mode to $8\times10^{-11}$, $3\times10^{-8}$, $3\times10^{-7}$; some real bulk modes never converge ($10^{-3}$--$10^{-2}$).
\item \textbf{Open chain, $L = 6$, scar initial states} (mixtures of spin coherent states $\dyad{\mathbf n}{\mathbf n}^{\otimes L}$, all inside the tower) against random ones, dense, $T = 200$: the tower modes converge 2--4 orders of magnitude better ($\lambda = 0$: $8\times10^{-11}$ against $2\times10^{-9}$; $-0.02 + i$: $3\times10^{-9}$ against $10^{-5}$; $-0.04 + 2i$: $3\times10^{-7}$ against $5\times10^{-3}$).
\item \textbf{Truncation} (open chain, scar initial states, $\chi = 16$ and $32$ of the full $64$): only $\lambda = 0$ ($4\times10^{-8}$, $2\times10^{-7}$) and the first edge mode ($9\times10^{-6}$, $2\times10^{-5}$) converge; every other tower mode stays at $\approx10^{-3}$, and the trace of the traceless state leaks to $3\times10^{-3}$ and $7\times10^{-3}$. This is not a time-step error: at $t \le 10$ and $\chi = 32$ the Krylov matrices differ from the dense ones by $3\times10^{-3}$ for internal steps $0.25$ and $0.05$ alike. So $\rho(t)$ from scar states needs nearly the full bond dimension already at $L = 6$: the weight that the dephasing moves out of the tower, at rate $\approx\gamma(L-1)$ (Eq.~\eqref{eq:cap}), is thermalized by $H_I$ and carries large operator entanglement.
\end{itemize}


\subsection{Forward/backward evolution with matrix product states}
\label{sec:fb}

\texttt{julia/forward\_backward.jl} implements Eq.~\eqref{eq:meet} (Sec.~\ref{sec:meet}) in Pauli coordinates: $\rho$ is evolved with $G$, the observable coordinates $o_P = \Tr(AP)/2^L$ with $G^{\mathsf T}$, and $\langle A\rangle_{t_1+t_2} = o(t_2)\cdot r(t_1)/\Tr\rho(t_1)$. $G^{\mathsf T}$ is built as an MPO from the same terms: the Hamiltonian part with $H \to -H$, the dissipator diagonal unchanged, the feed of decay and thermal noise transposed ($\mathds 1 \to \sigma^z$); checked against the transpose of the contracted $G$ to $5\times10^{-15}$ at $L = 6$. The same run also gives the forward-only and backward-only values at the same $\chi$, and the spread of $\langle A\rangle_t$ over all splits $t_1 + t_2 = t$, a truncation-error estimate.

\begin{itemize}
\item \textbf{Time stepping.} Two-site TDVP is wrong while the bond dimension is small: one step $\delta t = 0.1$ from $\ket{-}^{\otimes 6}$ is off by $10^{-2}$ in norm, with or without global Krylov expansion, scaling as $\delta t^2$ (the range-4 terms of $H$ leave the two-site tangent space). A Taylor series of $e^{G\tau}$ with exact MPO applications has no such error ($6\times10^{-8}$ for the same step, $\chi$ stays 1 as it must at $\gamma = 0$) but costs $\approx 25$\,s per substep at $\chi = 64$. The code uses Taylor while $\chi < 16$ and TDVP above.
\item \textbf{Validation at $L = 6$} (full $\chi = 64$, $\delta t = 0.1$, $t \le 5$, against the exact Rust evolution). Scar at $\gamma = 0$ (perfect revivals): largest deviation of $m_x$, $m_y$ $2.5\times10^{-6}$ (forward only), $1.3\times10^{-6}$ (meet in the middle), $8\times10^{-9}$ (backward only); all six magnetizations below $3\times10^{-6}$. Scar, $\sigma^z$ dephasing, $\gamma = 0.1$: below $3\times10^{-7}$ (forward, middle) and $5\times10^{-9}$ (backward). Other channels at $\gamma = 0.1$: scar under $\sigma^x$ dephasing $9\times10^{-6}$ and under decay $2\times10^{-6}$ (forward, middle; backward $\le 7\times10^{-9}$); $x$-N\'eel ($s_x$, $s_y$) under $\sigma^z$ dephasing and decay, and $z$-N\'eel ($s_z$) under $\sigma^z$, $\sigma^x$ dephasing and decay: $\le 6\times10^{-8}$ for all three. The backward half is the more accurate one from the scar state, because the forward state stays at small $\chi$ (where TDVP is not used, but the Taylor steps are compressed with a cutoff) for longer.
\item \textbf{Cost.} At $L = 6$ and full $\chi$ one run ($t \le 5$, one or two observables) takes 150--260\,s on the laptop (8 threads); with all six observables 370\,s ($\gamma = 0.1$) and 1070\,s ($\gamma = 0$). Dense vectors are far cheaper at this size; the method is meant for $L \gtrsim 14$, where it has not been run.
\end{itemize}

\subsection{Eigenvalues without LAPACK}
\label{sec:nolapack}
The entanglement spectra are computed in \texttt{rust/src/entanglement.rs} without external libraries, so the same binaries run on every machine. A real symmetric $n\times n$ matrix is brought to tridiagonal form by $n-2$ Householder reflections $H_k = \mathds 1 - \beta vv^{\mathsf T}$ acting on the trailing block, $B \to B - vw^{\mathsf T} - wv^{\mathsf T}$ with $p = \beta Bv$, $w = p - \tfrac\beta2(v^{\mathsf T}p)\,v$. The column is first divided by its largest entry: a tiny column would otherwise make $v^{\mathsf T}v$ underflow and $\beta = 2/v^{\mathsf T}v$ overflow (this produced NaN for the scar state, whose amplitudes carry round-off imaginary parts $\approx 4\times10^{-17}$). The tridiagonal eigenvalues follow from implicit QL iterations with Wilkinson shifts; an off-diagonal element is set to zero when it is below $\varepsilon$ times the sum of its two diagonal neighbours \emph{or} below $\varepsilon$ times the largest row sum of the tridiagonal matrix (the absolute test is needed for the many zero eigenvalues of a low-rank $\rho_A$). A Hermitian $M = X + iY$ is handled through the real symmetric matrix $\begin{pmatrix}X & -Y\\ Y & X\end{pmatrix}$, which has every eigenvalue of $M$ twice. Speed on one core $\approx 5$ GFLOP/s ($n = 1024$ in $0.3$ s, $n = 2048$ in $3.4$ s); inside the parallel loops over trajectories or single-flip runs it runs on one thread per matrix.

\paragraph{Reachable sizes, against LAPACK.} The solver returns eigenvalues only, of real symmetric or Hermitian matrices, at $O(n^3)$ (an $n = 4096$ Hermitian matrix, i.e.\ $8192$ in the real embedding, takes $\approx 3.5$ minutes on one core by extrapolation from the measured $n = 2048$). That is enough for every entanglement spectrum computed here, where the size is set by the state, not by the solver: pure-state block entanglement up to $L = 24$ ($\rho_A$ of dimension $2^{12} = 4096$) and operator entanglement of $\rho$ up to $L = 12$ ($RR^{\mathsf T}$ of dimension $4^6 = 4096$, $\approx 10$\,s per evaluation with 16 threads). Everything that needs eigenvectors or non-symmetric matrices uses LAPACK (Intel MKL): all eigenstates of $H$ (\texttt{zheevd}, $L \le 12$ here; $L = 14$, $n = 16384$, is the practical limit), the Liouvillian (\texttt{dgeev} on the dense real $4^L \times 4^L$ generator, $L \le 7$: $4.4$ GB and $18$ minutes at $L = 7$, $\approx 100$ GB at $L = 8$), and the generalized problems for degenerate modes (\texttt{zhegv}).
